\documentclass[a5paper,10pt]{article}

% https://github.com/InDevRus/filippov-solutions

% Нумерация страниц
\pagenumbering{gobble}

% Настройка полей
\usepackage{geometry}
\geometry{left=1cm}
\geometry{right=1cm}
\geometry{top=1cm}
\geometry{bottom=1.5cm}

% Вставка таблицы
\usepackage{array}
\usepackage{tabularx}

% Инструменты выравнивания формул
\usepackage{amsmath}
\usepackage{mathtools}

% Диакритика в математике
\usepackage{amsfonts}

% Вычисления размеров на ходу
\usepackage{calc}

% Удобные записи производных
\usepackage[italic=true]{derivative}

% Настройки сносок
\usepackage{enotez}

% Длинные знаки векторов
\usepackage{anyfontsize}
\usepackage[b]{esvect}

% Вставка картинок
\usepackage{graphicx}

% Вставка красной строки
\usepackage{indentfirst}
\setlength{\parindent}{1em}

% Условия в командах
\usepackage{xifthen}

% Луа-скрипты
\usepackage{luacode}

% Настройка команд отдельно в математическом и текстовом режимах
\usepackage{mathcommand}

% Для повторения LaTeX кода
\usepackage{multido}

% Конфигурация отступов формул
\delimitershortfall-1sp
\usepackage{mleftright}
\mleftright

% Растягивание объектов
\usepackage{relsize}

% Обтекание текстом
\usepackage{wrapfig}
\usepackage{subcaption}
\usepackage{floatflt}

% Поддержка ввода кириллицы
\usepackage[english, russian]{babel}

% Настройка корневой позиции для компилятора
\usepackage{import}
\providecommand*{\ProjectRootPath}{..}

\import{\ProjectRootPath/preambles/macros/}{theorems}

\import{\ProjectRootPath/preambles/fonts}{font_setup}

% Знак градуса
\usepackage{siunitx}

\import{\ProjectRootPath/preambles/macros/}{operators}
\import{\ProjectRootPath/preambles/macros/}{redefinitions}

\import{\ProjectRootPath/preambles/fonts}{letters_setup}
\import{\ProjectRootPath/preambles/fonts/adjustments}{custom_superscripts}

\import{\ProjectRootPath/preambles/custom}{formatting}
\import{\ProjectRootPath/preambles/custom}{frames}
\import{\ProjectRootPath/preambles/custom}{list_setup}

% TODO: перевести команды на современный стиль объявления
\input{../preambles/plotting}

\begin{document}

$\ddot{x}\br{t} = 2x - 2x^3$.

$\tsys{& \dot{x}\br{t} = y \\& \dot{y}\br{t} = 2x - 2x^3}$;
$\dfrac {dy} {dx} = \dfrac {2x - 2x^3} {y}$.
$y^2 + x^4 - 2x^2 = C$.
$y^2 + \br{x^2 - 1}^2 = C$, $C > 0$. \vspace{1mm}

Исследуем особые точки $K\br{0; 0}$, $L\br{-1; 0}$ и $M\br{1; 0}$. 

\begin{enumerate}
	\item $K\br{0; 0}$. $\dfrac {dy} {dx} = \dfrac {2x - 2x^3} {y} = \dfrac {2x + O\br{x^3}} {y}$.
	$\begin{pmatrix} 0 & 1 \\ 2 & 0 \end{pmatrix}$.
	$\lambda^2 - 2 = 0$.
	$\lambda_{1} = -\sqrt{2}$.
	$\lambda_{2} = \sqrt{2}$.
	Особая точка $K\br{0; 0}$ -- <<седло>>.
	
	\item Точки $L\br{-1; 0}$ и $M\br{1; 0}$ имеют одинаковое поведение вокруг себя, так как траектории симметричны относительно оси $Oy$.
	
	Вместо замены переменных оптимальнее будет разложить числитель по формуле Тейлора в окрестности особой точки.
	
	Для первой имеем $\dfrac {dy} {dx} = \dfrac {-4\br{x + 1} + 6\br{x + 1}^2 - 2\br{x + 1}^3} {y} = \dfrac {-4\br{x + 1} + O\br{x + 1}^2} {y}$,
	так как первая, вторая и т. д. производные числителя равны соответственно $-4$, $12$, $-12$, $0$, $\ldots$
	
	Аналогично, для второй точки получаем
	$\dfrac {dy} {dx} = \dfrac {-4\br{x - 1} - 6\br{x - 1}^2 - 2\br{x - 1}^3} {y} = \linebreak =
	\dfrac {-4\br{x - 1} + O\br{x - 1}^2} {y}$.
	
	Для обеих матрица оператора равна $\begin{pmatrix} 0 & 1 \\ -4 & 0 \end{pmatrix}$, $\lambda^2 + 4 = 0$. Из-за обозначенной выше симметрии обе особые точки -- <<центры>>.
\end{enumerate}

На графике изображены траектории для разных $C > 0$.
Фиолетовые кривые соответствуют $C \in \br{0; 1}$, зелёные $C > 1$, а чёрная $C = 1$.

\begin{figure}[h]
\begin{minipage}{.5\textwidth}
	\centering
	\begin{tikzpicture}
		\pgfplotsset{
			Graph/.style = {
				thin,
				samples = 500,
				restrict y to domain = -5 : 5,
			},
			DirectionGraph/.style = {
				samples = 20,
				quiver = {scale arrows = 0.075},
				restrict y to domain = -5 : 5,
			},
			DirectionGraph1/.style = {
				DirectionGraph,
				-stealth,
				quiver = {
					u = {dot_x(x, y) / sqrt(dot_x(x, y) ^ 2 + dot_y(x, y) ^ 2)},
					v = {dot_y(x, y) / sqrt(dot_x(x, y) ^ 2 + dot_y(x, y) ^ 2)}
				}
			},
			DirectionGraph2/.style = {
				DirectionGraph,
				quiver = {
					u = {-dot_x(x, y) / sqrt(dot_x(x, y) ^ 2 + dot_y(x, y) ^ 2)},
					v = {-dot_y(x, y) / sqrt(dot_x(x, y) ^ 2 + dot_y(x, y) ^ 2)}
				}
			},
		}
		\begin{axis}[
			trig format plots = rad,
			width = 6.5cm,
			height = 10cm,
			ylabel=$\dot{x}$,
			ymin = -1.9,
			ymax = 1.9,
			xmin = -1.8,
			xmax = 1.8, 
			xtick = {-10, ..., 10},
			ytick = {-10, ..., 10},
			minor tick num = 3,
			declare function = {
				f(\C, \x) = sqrt(1 + sqrt(\C) * cos(x));
				g(\C, \x) = sqrt(\C) * sin(x);
				h(\C, \x) = sqrt(\C - (x^2 - 1)^2);
				dot_x(\x, \y) = y;
				dot_y(\x, \y) = 2 * x - 2 * x * x * x;
			}]
			
			\pgfplotsinvokeforeach {0.1, 0.2, 0.4, 0.6, 0.75, 0.9} {
				\foreach \s in {-1, 1} {
					\addplot[Graph, magenta, domain = -pi : pi] 
					({\s * f(#1, \x)}, {g(#1, \x)});
				}
			}
			
			\pgfplotsinvokeforeach {1} {
				\foreach \s in {-1, 1} {
					\addplot[Graph, black, domain = -pi : pi] 
					({\s * f(#1, \x)}, {g(#1, \x)});
				}
			}
			
			\pgfplotsinvokeforeach {1.2, 1.5, 1.8, 2.2} {
				\foreach \s in {-1, 1} {
					\addplot[Graph, olive, domain = -pi : pi] 
					({\s * f(#1, \x)}, {g(#1, \x)});
					\addplot[Graph, olive, domain = -0.1 : 0.1]
					{\s * h(#1, \x)};
				}
			}
		\end{axis}
	\end{tikzpicture}	
\end{minipage}
\begin{minipage}{.5\textwidth}
	\centering
	\begin{tikzpicture}
		\pgfplotsset{
			Graph/.style = {
				thin,
				samples = 500,
				restrict y to domain = -5 : 5,
			},
			DirectionGraph/.style = {
				samples = 20,
				quiver = {scale arrows = 0.075},
				restrict y to domain = -5 : 5,
			},
			DirectionGraph1/.style = {
				DirectionGraph,
				-stealth,
				quiver = {
					u = {dot_x(x, y) / sqrt(dot_x(x, y) ^ 2 + dot_y(x, y) ^ 2)},
					v = {dot_y(x, y) / sqrt(dot_x(x, y) ^ 2 + dot_y(x, y) ^ 2)}
				}
			},
			DirectionGraph2/.style = {
				DirectionGraph,
				quiver = {
					u = {-dot_x(x, y) / sqrt(dot_x(x, y) ^ 2 + dot_y(x, y) ^ 2)},
					v = {-dot_y(x, y) / sqrt(dot_x(x, y) ^ 2 + dot_y(x, y) ^ 2)}
				}
			},
		}
		\begin{axis}[
			trig format plots = rad,
			width = 6.5cm,
			height = 10cm,
			ylabel=$\dot{x}$,
			ymin = -1.9,
			ymax = 1.9,
			xmin = -1.8,
			xmax = 1.8, 
			xtick = {-10, ..., 10},
			ytick = {-10, ..., 10},
			minor tick num = 3,
			declare function = {
				f(\C, \x) = sqrt(1 + sqrt(\C) * cos(x));
				g(\C, \x) = sqrt(\C) * sin(x);
				h(\C, \x) = sqrt(\C - (x^2 - 1)^2);
				dot_x(\x, \y) = y;
				dot_y(\x, \y) = 2 * x - 2 * x * x * x;
			}]			
			\pgfplotsinvokeforeach {1} {
				\foreach \s in {-1, 1} {
					\addplot[Graph, black, domain = -pi : pi] 
					({\s * f(#1, \x)}, {g(#1, \x)});
					\addplot[DirectionGraph1, black, domain = -pi + 0.1 : pi - 0.1] 
					({\s * f(#1, \x)}, {g(#1, \x)});
				}
			}
		\end{axis}
	\end{tikzpicture}
\end{minipage}
\end{figure}

\pagebreak
Построенный фазовый портрет позволяет сделать вывод о поведении кривых при $t \to +\infty$. Пусть $x\br{t}$ -- решение задачи Коши $\ddot{x}\br{t} = 2x - 2x^3$, $x\br{t_0} = a$, $\dot{x}\br{t_0} = b$. Обозначим $b^2 + \br{a^2 + 1}^2 = C \ge 0$. При $C \ne 1$ любое решение $x\br{t}$ при $t \to +\infty$ принимает ограниченный повторяющийся набор значений, причём:
\begin{itemize}
	\item если $b^2 + \br{a^2 - 1}^2 = C \in \br{0; 1}$, то решения сохраняют знак и находятся в окрестности особой точки:
	\begin{itemize}
		\item при $a > 0$: $0 < \sqrt{1 - \sqrt{C}} \le x\br{t} \le \sqrt{1 + \sqrt{C}}$;
		\item при $a < 0$: $-\sqrt{1 + \sqrt{C}} \le x\br{t} \le -\sqrt{1 - \sqrt{C}} < 0$;
	\end{itemize} 
	\vspace{1mm}
	\item если $b^2 + \br{a^2 - 1}^2 = C > 1$, то решения знакопеременны, проходят через ноль бесконечное число раз, и
	$-\sqrt{1 + \sqrt{C}} \le x\br{t} \le \sqrt{1 + \sqrt{C}}$; \vspace{1mm}
	\item случаи $C = 0$, т. е.
	$\tsys{& a = -1 \\& b = 1}$ и $\tsys{& a = 1 \\& b = 1}$ в силу теоремы единственности соответствуют решениям $x\br{t} = -1$ и $x\br{t} = 1$ соответственно;
	\item случай $a = b = 0$ даёт нулевое решение $x\br{t} = 0$.
\end{itemize}
Особого интереса заслуживает случай $b^2 + \br{a^2 - 1}^2 = C = 1$.
Каждая такая задача Коши имеет на фазовой плоскости такую траекторию, которая не может пересекать начало координат, иначе при конечном $t$ решение стало бы тождественно равным нулю. Следовательно, решение такой задачи Коши сохраняет знак и стремится к нулю при $t \to +\infty$.
В частности, иллюстрациями к этому случаю будут две задачи Коши со следующими решениями
\begin{align*}
	\tsys{& \ddot{x}\br{t} = 2x - 2x^3 \\& x\br{t_{0}} = 0 \\& \dot{x}\br{t_0} = \pm \sqrt{2}} 
	&& x\br{t} = \pm \dfrac {2} {\sqrt{2\sqrt{2} \ch\br{t - t_{0}} + 1}}.
\end{align*}
\begin{figure}[!h]
	\centering
	\begin{tikzpicture}
		\pgfplotsset{
			Graph/.style = {
				samples = 500,
				restrict y to domain = -5 : 5,
			},
		}
		\begin{axis}[
			trig format plots = rad,
			xlabel=$t$,
			ylabel=$x$,
			width = 14cm,
			height = 6.5cm,
			ymin = -1.9,
			ymax = 1.9,
			xmin = -4.2,
			xmax = 4.2, 
			xtick = {-10, ..., 10},
			ytick = {-10, ..., 10},
			minor tick num = 3,
			declare function = {
				f(\x) = 2 / sqrt(cosh(2 * sqrt(2) * x) + 1);
			}]			
			
			\pgfplotsinvokeforeach {1} {
				\foreach \s in {-1, 1} {
					\addplot[Graph, black, domain = -4.5 : 4.5] {\s * f(x)};
				}
			}
		\end{axis}
	\end{tikzpicture}
\end{figure}
\end{document}
